\documentclass{article}
\usepackage[UTF8]{ctex}
\usepackage{amsmath,amssymb,amsfonts}
\usepackage{enumitem}
\usepackage{geometry}
\geometry{a4paper, margin=1in}
% === 错题本排版优化（不参与题目标识） ===
\setlength{\parindent}{0pt}
\setlength{\parskip}{0.35em plus 0.10em minus 0.08em}
\setlength{\emergencystretch}{3em}
\setlist{itemsep=0.25em, topsep=0.40em, parsep=0pt}
\allowdisplaybreaks[2]
% === 错题本排版优化结束 ===
\begin{document}
\section*{苏州大学2026年数学分析真题}
\section*{90．苏州大学2026年数学分析试卷}

\subsection{题目1}
1．（20 分）解答如下问题：

\begin{itemize}
  \item[（1）] 计算极限 \(\lim _{n \rightarrow \infty} \lim _{x \rightarrow 0} \frac{\sin (n x)-n \sin x}{n^{3} x^{3}}\) ．
  \item[（2）] 设函数 \(f(x)=\left\{\begin{array}{ll}(1+x)^{\frac{1}{x}}, & x \neq 0, \\ e, & x=0 .\end{array}\right.\) 求 \(f^{\prime}(0)\) ．
\end{itemize}

\subsection{题目2}
2．（20 分）已知函数 \(f\) 在 \([0,1]\) 上连续，在 \((0,1)\) 上可微，且 \(f(0)=f(1)=0, f\left(\frac{1}{2}\right)=1\) ．

\begin{itemize}
  \item[（1）] 证明：存在 \(\eta \in\left(\frac{1}{2}, 1\right)\) ，使得 \(f(\eta)=\eta\) 。
  \item[（2）] 证明：存在 \(\xi \in(0,1)\) ，使得 \(f^{\prime}(\xi)-\lambda[f(\xi)-\xi]=1\) ．
\end{itemize}

\subsection{题目3}
3．（15 分）证明：在 \(\left[0, \frac{\pi}{2}\right]\) 上，有 \(\sin x \geq \frac{2}{\pi} x+\frac{1}{12 \pi} x\left(\pi^{2}-4 x^{2}\right)\) ．

\subsection{题目4}
4．（15 分）解答如下问题：

\begin{itemize}
  \item[（1）] 设 \(f(x)\) 连续，证明： \(\int_{n \rightarrow \infty} \int_{0}^{\pi} f(x)|\sin n x| \mathrm{d} x=\frac{2}{\pi} \int_{0}^{\pi} f(x) \mathrm{d} x\) ．
  \item[（2）] 设 \(f(x)\) 连续，证明： \(\lim _{n \rightarrow \infty}\left(\int_{a}^{b}|f(x)|^{n} \mathrm{~d} x\right)^{\frac{1}{n}}=\max _{x \in[a, b]}|f(x)|\) ．
\end{itemize}

\subsection{题目5}
5．（15 分）解答如下问题：

\begin{itemize}
  \item[（1）] 判断函数列 \(f_{n}(x)=n x e^{-n x^{2}}\) 在 \((-\infty,+\infty)\) 上的一致收敛性．
  \item[（2）] 判断函数项级数 \(\sum_{n=1}^{\infty} \frac{\ln \left(1+n^{2} x^{2}\right)}{n^{2}}\) 在 \((-2,2)\) 上的一致收敛性．
\end{itemize}

\subsection{题目6}
6．（20 分）解答如下问题：

\begin{itemize}
  \item[（1）] 求积分 \(\int_{0}^{1}(\ln x)^{2026} \mathrm{~d} x\) ．
  \item[（2）] 求积分 \(\int_{0}^{\frac{\pi}{2}} \ln (\sin x) \mathrm{d} x\) ．
\end{itemize}

\subsection{题目7}
7．（15 分）解答如下问题：

\begin{itemize}
  \item[（1）] 求 \(\int_{L}\left(x^{2}-y\right) \mathrm{d} x+\left(y^{2}-x\right) \mathrm{d} y\) ，其中 \(L\) 是 \(x^{2}+y^{2}=4\) 的上半部分，方向从（－2，0）到（2，0）。
  \item[（2）] 求由曲面 \(z=\sqrt{2-x^{2}-y^{2}}\) 与抛物面 \(z=x^{2}+y^{2}\) 所围立体部分的表面积。
\end{itemize}

\subsection{题目8}
8．（15 分）设 \(f(x, y, z)=2\left(x^{3}+y^{3}+z^{3}\right)-5\left(x^{2}+y^{2}+z^{2}\right)\) ，其中 \(x, y, z \geq 0\) ，求 \(f\) 在条件 \(x+y+z=3\)下的最值。

\subsection{题目9}
9．（15 分）二元函数 \(u(x, y)\) 在区域 \(D=[0,1] \times[0,1]\) 上具有二阶连续偏导数，且在边界 \(\partial D\) 上 \(u=0\) ．如果

\[
\frac{\partial^{2} u}{\partial x^{2}}(x, y)+\frac{\partial^{2} u}{\partial y^{2}}(x, y)=-\lambda u(x, y), \forall(x, y) \in D .
\]

\begin{itemize}
  \item[（1）] 证明： \(\iint_{D}\left[\left(\frac{\partial u}{\partial x}\right)^{2}+\left(\frac{\partial u}{\partial y}\right)^{2}\right] \mathrm{d} x \mathrm{~d} y=\lambda \iint_{D} u^{2} \mathrm{~d} x \mathrm{~d} y\) 。
  \item[（2）] 若存在函数 \(f, g\) ，使得 \(u\) 可以表示为 \(u(x, y)=f(x) g(y)\) ，且在区域 \(D\) 内部 \(u>0\) ，求 \(\lambda\) 的值．
\end{itemize}
\end{document}
